\documentclass[10pt,a4paper]{amsart}
\usepackage[margin=19mm]{geometry}
\usepackage{amsmath,amssymb,mathtools}
\usepackage[scr=rsfs,cal=euler]{mathalfa}
\usepackage{xfrac}
\usepackage[unicode,hidelinks,bookmarks=false]{hyperref}
\newcommand{\ie}{i.e.\ }
\newcommand{\cf}{cf.\ }
\newcommand{\resp}{respectively\ }
\newcommand{\Subsection}{\subsection}
\providecommand{\nobreakdash}{\nobreak\hbox{-}}
\setlength{\emergencystretch}{2em}
\multlinegap0pt
\usepackage{siunitx}
\usepackage{siunitx}
\usepackage{siunitx}
\usepackage{siunitx}
\usepackage{bm}
\usepackage{amssymb}
\usepackage{float}
\usepackage{mathtools}
\usepackage{dsfont}
\usepackage{siunitx}
\usepackage{tikz}
\usepackage[shortlabels]{enumitem}
\usepackage{algorithm}\usepackage{algpseudocode}
\newtheorem{lemma}{Lemma}
\newcommand{\bQ}{\bm Q}
\newcommand{\bSigma}{\bm{\Sigma}}
\newcommand{\bmean}{\bm{m}}
\newcommand{\bs}{\bm{s}}
\newcommand{\bz}{\bm{z}}
\newcommand{\mN}{\mathcal{N}}
\newcommand{\mbGM}{\mathbb{GM}}
\newcommand{\mbR}{\mathbb R}
\newcommand{\mbS}{\mathbb S}
\title{Robust Chance-Constrained Optimization using a Continuous Parameter Space Wasserstein-2 Ambiguity Set of Gaussian Mixtures}
\author{Shibshankar Dey, Sanjay Mehrotra}
\date{Source: arXiv:2607.17018v2 (CC BY 4.0)}
\begin{document}
\maketitle

\noindent\emph{Source and licence.} Robust Chance-Constrained Optimization using a Continuous Parameter Space Wasserstein-2 Ambiguity Set of Gaussian Mixtures. Version arXiv:2607.17018v2 (CC BY 4.0), distributed by arXiv under the Creative Commons Attribution 4.0 International licence, http://creativecommons.org/licenses/by/4.0/, as recorded in the arXiv licence metadata for this exact version and shown on the abstract page https://arxiv.org/abs/2607.17018v2. Full licence notice: \texttt{licenses/arxiv-2607.17018-CC-BY-4.0.txt}.\par\smallskip

\noindent\textit{Editorial scope.} Counted unit: the article's lemma lem:nonempty-restricted-GM-coupling (source0.tex lines 4043--4054). The article's complete original proof of it is delivered with it (source0.tex lines 4056--4142, one \texttt{proof} environment of 87 lines). No mathematics is written, repaired or re-derived: the statement and the proof are the article's own lines, in the article's own notation, and no retained line is substituted or re-flowed.  No further source-local context is needed: every cross-reference of every retained line resolves inside the two delivered spans.  Nothing was omitted from any retained span, so the package carries no omission record.  No retained line contains a citation, so the wrapper carries no bibliography and no bibliography entry.  No retained line contains a figure environment, an image-inclusion command or drawing code, so no image is delivered and the package declares no asset dependency.  Editorial formatting only, in addition: an AMS \texttt{amsart} preamble that restates the article's own declarations and macros used by the retained lines, namely the environments \texttt{lemma} on separate counters, together with the standard packages those lines need.  The retained environments print in this document as Lemma 1 in order of appearance, whereas the article prints the same items with the numbering of its own full document.  The complete native source of the article as distributed in the arXiv e-print for this exact version is bundled unchanged as \texttt{sources/205523/source0.tex}.
\begin{lemma}[Nonemptiness under restricted Gaussian-mixture support]
\label{lem:nonempty-restricted-GM-coupling}
Suppose that \(\mu_1,\mu_2\in\mathcal P(\mbR^n)\) admit mixing laws over
\(\mathcal S\); that is, there exist \(\eta_1,\eta_2\in\mathcal P(\mathcal S)\)
such that, for every \(C\in\mathcal B(\mbR^n)\),
\[
\mu_1(C)=\int_{\mathcal S}\mN(\bs)(C)\,\eta_1(d\bs),
\qquad
\mu_2(C)=\int_{\mathcal S}\mN(\bs)(C)\,\eta_2(d\bs).
\]
Then $\quad \Pi(\mu_1,\mu_2)\cap \mbGM_{2n}^{\mathcal S}\neq\varnothing.$
\end{lemma}

\begin{proof}
Let\vspace{-2em}
\[
\pi:=\eta_1\otimes\eta_2\in\mathcal P(\mathcal S\times\mathcal S).
\]
For \(\bs_i=(\bmean_i,\bQ_i)\in\mathcal S\), \(i=1,2\), let
\(\mathcal G(\bs_1,\bs_2)\) denote the Gelbrich coupling between
\(\mN(\bs_1)\) and \(\mN(\bs_2)\). Thus
\[
\mathcal G(\bs_1,\bs_2)
=
\mN\!\left(
(\bmean_1,\bmean_2),
\overline{\bQ}(\bs_1,\bs_2)
\right),
\]
where\vspace{-1.75em}
\[
\overline{\bQ}(\bs_1,\bs_2)
=
\begin{pmatrix}
\bQ_1 & \bSigma(\bQ_1,\bQ_2)\\
\boldsymbol\Sigma(\bQ_1,\bQ_2)^\top & \bQ_2
\end{pmatrix}
\in\mbS^{2n}_+.
\]
Since the diagonal covariance blocks of \(\overline{\bQ}(\bs_1,\bs_2)\) are
\(\bQ_1\) and \(\bQ_2\), the map $\mathcal T(\bs_1,\bs_2)
:=
\left(
(\bmean_1,\bmean_2),
\overline{\bQ}(\bs_1,\bs_2)
\right)$ satisfies $\mathcal T(\bs_1,\bs_2)\in\mathcal Z_{\mathcal S}
\quad
\forall (\bs_1,\bs_2)\in\mathcal S\times\mathcal S$. Hence, $\lambda:=\mathcal T_{\#}\pi\in\mathcal P(\mathcal Z_{\mathcal S})$. 

Define \(\nu\in\mathcal P(\mbR^{2n})\) by
\[
\nu(A):=\int_{\mathcal Z_{\mathcal S}}\mN(\bz)(A)\,\lambda(d\bz),
\qquad
A\in\mathcal B(\mbR^{2n}).
\]
Equivalently,
\[
\nu(A)
=
\int_{\mathcal S\times\mathcal S}
\mathcal G(\bs_1,\bs_2)(A)\,\pi(ds_1,ds_2).
\]
Therefore, \(\nu\in\mbGM_{2n}^{\mathcal S}\).

It remains to show that \(\nu\in\Pi(\mu_1,\mu_2)\). Let
\(B\in\mathcal B(\mbR^n)\). Since \(\mathcal G(\bs_1,\bs_2)\) is a coupling of
\(\mN(\bs_1)\) and \(\mN(\bs_2)\),
\[
\mathcal G(\bs_1,\bs_2)(B\times\mbR^n)
=
\mN(\bs_1)(B).
\]
Thus\vspace{-1.85em}
\[
\nu(B\times\mbR^n)
=
\int_{\mathcal S\times\mathcal S}
\mN(\bs_1)(B)\,(\eta_1\otimes\eta_2)(ds_1,ds_2)
=
\int_{\mathcal S}\mN(\bs_1)(B)\,\eta_1(ds_1)
=
\mu_1(B).
\]
Hence, \(\nu^{(1)}=\mu_1\). The same argument gives
\[
\nu(\mbR^n\times B)
=
\int_{\mathcal S}\mN(\bs_2)(B)\,\eta_2(ds_2)
=
\mu_2(B),
\]
so \(\nu^{(2)}=\mu_2\). Therefore, \(\nu\in\Pi(\mu_1,\mu_2)\).

Combining \(\nu\in\Pi(\mu_1,\mu_2)\) and
\(\nu\in\mbGM_{2n}^{\mathcal S}\), we obtain
\[
\nu\in\Pi(\mu_1,\mu_2)\cap\mbGM_{2n}^{\mathcal S}.
\]
Hence, the intersection is nonempty.
\end{proof}

\end{document}
